% =====================================================================
%  Phương pháp — GAS (Generative Aspect Sentiment)
%  Baseline sinh văn bản cho Stage 2 (ATE); chỉ đánh giá aspect term.
%  Dùng % =====================================================================
%  Phương pháp — GAS (Generative Aspect Sentiment)
%  Baseline sinh văn bản cho Stage 2 (ATE); chỉ đánh giá aspect term.
%  Dùng % =====================================================================
%  Phương pháp — GAS (Generative Aspect Sentiment)
%  Baseline sinh văn bản cho Stage 2 (ATE); chỉ đánh giá aspect term.
%  Dùng % =====================================================================
%  Phương pháp — GAS (Generative Aspect Sentiment)
%  Baseline sinh văn bản cho Stage 2 (ATE); chỉ đánh giá aspect term.
%  Dùng \input{phuong_phap_gas} trong main.tex
% =====================================================================

% =====================================================================
\section{Phương pháp GAS (Generative Aspect Sentiment)}
\label{sec:gas}
% =====================================================================

GAS (\textit{Generative Aspect Sentiment}) là phương pháp trích xuất
ABSA một giai đoạn được đề xuất bởi Zhang et al.~\cite{zhang2021gas}.
Khác với các tiếp cận phân tầng truyền thống — đầu tiên trích xuất
aspect term (ATE), sau đó phân loại cảm xúc (APC) — GAS sử dụng một
mô hình sinh duy nhất để dự đoán đồng thời ba nhãn trong một lần giải
mã:
\textbf{aspect term}, \textbf{aspect category}, và
\textbf{sentiment polarity}.

Trong luận văn này, GAS được sử dụng làm \textit{baseline cho bước
trích xuất aspect term} (Stage~2 — ATE).  Mô hình vẫn được huấn luyện
để sinh bộ ba $(a_i, c_i, s_i)$ theo đúng cấu trúc gốc, song kết quả
so sánh chỉ tập trung vào nhãn \textbf{aspect term}: category và
sentiment không thuộc phạm vi đánh giá của chương này.

% ---------------------------------------------------------------------
\subsection{Kiến trúc tổng quan}
\label{subsec:gas_arch}
% ---------------------------------------------------------------------

GAS xây dựng bài toán ABSA như một bài toán dịch chuỗi
(\textit{sequence-to-sequence}), với mô hình nền là
\texttt{T5ForConditionalGeneration}~\cite{raffel2020t5} phiên bản
\texttt{t5-base}.  Với mỗi câu đầu vào $\mathbf{x}$, mô hình được
huấn luyện để sinh ra chuỗi đích $\mathbf{y}$ theo cú pháp cố định:

\begin{equation}
  \mathbf{y} = \texttt{(}a_1,\, c_1,\, s_1\texttt{);\;
               (}a_2,\, c_2,\, s_2\texttt{);\;} \dots
  \label{eq:gas_format}
\end{equation}

\noindent trong đó $a_i$ là cụm từ khía cạnh (\textit{aspect term}),
$c_i \in \{\text{SERVICE, FACILITY, AMENITY, EXPERIENCE, BRANDING,
LOYALTY}\}$ là danh mục khía cạnh, và
$s_i \in \{\text{positive, negative, neutral}\}$ là cực tính cảm xúc.
Nếu câu không chứa khía cạnh nào, mô hình sinh ra chuỗi
\texttt{"none"}.  Trong luận văn này chỉ trường $a_i$ được sử dụng
cho mục đích so sánh; $c_i$ và $s_i$ được Stage~3 (APC) xử lý riêng.

Quá trình giải mã sử dụng \textit{beam search} với chiều rộng tia
$k = 4$ và điều kiện dừng sớm (\textit{early stopping}) khi toàn bộ
các tia đều kết thúc bằng token \texttt{</s>}.

% ---------------------------------------------------------------------
\subsection{Tiền xử lý dữ liệu}
\label{subsec:gas_data}
% ---------------------------------------------------------------------

Dữ liệu đầu vào được đọc từ định dạng \texttt{.apc} bốn dòng lặp lại
(câu — aspect term — aspect category — sentiment).
Vì định dạng này lưu một dòng cho mỗi cặp \textit{(câu, aspect)},
các mẫu cùng thuộc một câu được nhóm lại thành một
\textit{record} duy nhất.  Chuỗi đích tương ứng được xây dựng bằng
cách nối tất cả triplet của câu đó theo cú pháp~\eqref{eq:gas_format}.

Mỗi record có dạng:

\begin{equation}
  r = \bigl(\underbrace{x}_{\text{input\_text}},\;
            \underbrace{y}_{\text{target\_text}},\;
            \underbrace{\{a_i\}}_{\text{aspects}},\;
            \underbrace{\{(a_i, c_i, s_i)\}}_{\text{gold\_triples}}
      \bigr)
  \label{eq:gas_record}
\end{equation}

Tokenization được thực hiện bởi \texttt{T5Tokenizer} với
\texttt{max\_length}$= 128$ cho cả đầu vào lẫn đầu ra.  Các vị trí
padding trong chuỗi nhãn được thay bằng $-100$ để hàm mất mát
cross-entropy bỏ qua.

% ---------------------------------------------------------------------
\subsection{Huấn luyện}
\label{subsec:gas_train}
% ---------------------------------------------------------------------

Mô hình được tinh chỉnh bằng hàm mất mát cross-entropy chuẩn của
seq2seq (T5 tự tính theo cơ chế \textit{teacher forcing}):

\begin{equation}
  \mathcal{L}_{\text{GAS}} =
    -\sum_{t=1}^{T} \log P\!\left(y_t \mid y_{<t},\, \mathbf{x};\,
    \Theta\right)
  \label{eq:gas_loss}
\end{equation}

\noindent với $T$ là độ dài chuỗi đích và $\Theta$ là toàn bộ tham số
mô hình.  Quá trình tối ưu sử dụng AdamW với learning rate
$3 \times 10^{-4}$, kết hợp \textit{linear warmup} trong 10\% tổng
số bước.  Kỹ thuật mixed precision (AMP) được bật khi có GPU CUDA.

Early stopping theo dõi \textbf{ATE F1} (Aspect Term F1) trên tập
phát triển (patience$= 4$ epoch): một dự đoán được tính là True
Positive khi aspect term dự đoán khớp chính xác với nhãn vàng, bất
kể category hay sentiment — nhất quán với mục tiêu của Stage~2 trong
pipeline.

% ---------------------------------------------------------------------
\subsection{Hậu xử lý kết quả sinh}
\label{subsec:gas_postprocess}
% ---------------------------------------------------------------------

Chuỗi văn bản do T5 sinh ra được phân tích cú pháp bằng biểu thức
chính quy để trích xuất danh sách triplet:

\begin{equation}
  \texttt{\_GAS\_RE} =
    \texttt{\textbackslash(\textbackslash s*(.+?),\textbackslash s*%
    (.+?),\textbackslash s*(positive|negative|neutral)\textbackslash s*%
    \textbackslash)}
  \label{eq:gas_regex}
\end{equation}

Group~1, 2, 3 lần lượt trả về $a_i$, $c_i$, $s_i$.  Nếu chuỗi sinh
là \texttt{"none"} hoặc không khớp mẫu nào, hàm trả về danh sách
rỗng.

\paragraph{Chuẩn hóa aspect term bằng Levenshtein.}
Mô hình sinh đôi khi tạo ra cụm từ khía cạnh với lỗi chính tả nhỏ
(thiếu/thừa ký tự, hoán vị).  Để đảm bảo aspect term kết quả
\textit{thực sự xuất hiện trong câu gốc}, bước chuẩn hóa n-gram
Levenshtein được áp dụng sau khi parse:

\begin{enumerate}
  \item \textbf{Xây dựng từ điển n-gram:}  Từ câu gốc $\mathbf{x}$,
    tập ứng viên $\mathcal{V}$ gồm tất cả chuỗi con liên tục theo từ
    (unigram, bigram, trigram, \dots).

  \item \textbf{Ánh xạ về n-gram gần nhất:}  Nếu aspect term dự đoán
    $\hat{a}$ chưa thuộc $\mathcal{V}$, thay bằng:
    \begin{equation}
      a^* = \arg\min_{v \in \mathcal{V}}\;
        \mathrm{Lev}(\hat{a},\, v)
      \label{eq:levenshtein}
    \end{equation}
    \noindent trong đó $\mathrm{Lev}(\cdot, \cdot)$ là khoảng cách
    chỉnh sửa Levenshtein.  Nếu $\hat{a} \in \mathcal{V}$, giữ
    nguyên.
\end{enumerate}

Chỉ aspect term được chuẩn hóa; category và sentiment giữ nguyên
như mô hình sinh ra.

% ---------------------------------------------------------------------
\subsection{Độ đo đánh giá}
\label{subsec:gas_eval}
% ---------------------------------------------------------------------

Vì luận văn chỉ sử dụng GAS cho bước trích xuất aspect term (Stage~2),
độ đo duy nhất được báo cáo là \textbf{ATE F1} (Aspect Term F1):
một dự đoán được tính là True Positive khi chuỗi aspect term dự đoán
khớp chính xác (\textit{exact match}) với một aspect term trong nhãn vàng
của câu — category và sentiment không tham gia so khớp.

Precision, Recall và F1 được tính theo công thức chuẩn:

\begin{equation}
  P = \frac{\mathrm{TP}}{\mathrm{TP} + \mathrm{FP}}, \quad
  R = \frac{\mathrm{TP}}{\mathrm{TP} + \mathrm{FN}}, \quad
  F_1 = \frac{2PR}{P + R}
  \label{eq:prf}
\end{equation}

\noindent trong đó TP, FP, FN được tổng hợp (\textit{micro}) qua toàn
bộ câu của tập kiểm thử trước khi tính tỉ lệ — nhất quán với giao thức
đánh giá ATE trong các nghiên cứu ABSA~\cite{zhang2021gas,luo2019doer}.
 trong main.tex
% =====================================================================

% =====================================================================
\section{Phương pháp GAS (Generative Aspect Sentiment)}
\label{sec:gas}
% =====================================================================

GAS (\textit{Generative Aspect Sentiment}) là phương pháp trích xuất
ABSA một giai đoạn được đề xuất bởi Zhang et al.~\cite{zhang2021gas}.
Khác với các tiếp cận phân tầng truyền thống — đầu tiên trích xuất
aspect term (ATE), sau đó phân loại cảm xúc (APC) — GAS sử dụng một
mô hình sinh duy nhất để dự đoán đồng thời ba nhãn trong một lần giải
mã:
\textbf{aspect term}, \textbf{aspect category}, và
\textbf{sentiment polarity}.

Trong luận văn này, GAS được sử dụng làm \textit{baseline cho bước
trích xuất aspect term} (Stage~2 — ATE).  Mô hình vẫn được huấn luyện
để sinh bộ ba $(a_i, c_i, s_i)$ theo đúng cấu trúc gốc, song kết quả
so sánh chỉ tập trung vào nhãn \textbf{aspect term}: category và
sentiment không thuộc phạm vi đánh giá của chương này.

% ---------------------------------------------------------------------
\subsection{Kiến trúc tổng quan}
\label{subsec:gas_arch}
% ---------------------------------------------------------------------

GAS xây dựng bài toán ABSA như một bài toán dịch chuỗi
(\textit{sequence-to-sequence}), với mô hình nền là
\texttt{T5ForConditionalGeneration}~\cite{raffel2020t5} phiên bản
\texttt{t5-base}.  Với mỗi câu đầu vào $\mathbf{x}$, mô hình được
huấn luyện để sinh ra chuỗi đích $\mathbf{y}$ theo cú pháp cố định:

\begin{equation}
  \mathbf{y} = \texttt{(}a_1,\, c_1,\, s_1\texttt{);\;
               (}a_2,\, c_2,\, s_2\texttt{);\;} \dots
  \label{eq:gas_format}
\end{equation}

\noindent trong đó $a_i$ là cụm từ khía cạnh (\textit{aspect term}),
$c_i \in \{\text{SERVICE, FACILITY, AMENITY, EXPERIENCE, BRANDING,
LOYALTY}\}$ là danh mục khía cạnh, và
$s_i \in \{\text{positive, negative, neutral}\}$ là cực tính cảm xúc.
Nếu câu không chứa khía cạnh nào, mô hình sinh ra chuỗi
\texttt{"none"}.  Trong luận văn này chỉ trường $a_i$ được sử dụng
cho mục đích so sánh; $c_i$ và $s_i$ được Stage~3 (APC) xử lý riêng.

Quá trình giải mã sử dụng \textit{beam search} với chiều rộng tia
$k = 4$ và điều kiện dừng sớm (\textit{early stopping}) khi toàn bộ
các tia đều kết thúc bằng token \texttt{</s>}.

% ---------------------------------------------------------------------
\subsection{Tiền xử lý dữ liệu}
\label{subsec:gas_data}
% ---------------------------------------------------------------------

Dữ liệu đầu vào được đọc từ định dạng \texttt{.apc} bốn dòng lặp lại
(câu — aspect term — aspect category — sentiment).
Vì định dạng này lưu một dòng cho mỗi cặp \textit{(câu, aspect)},
các mẫu cùng thuộc một câu được nhóm lại thành một
\textit{record} duy nhất.  Chuỗi đích tương ứng được xây dựng bằng
cách nối tất cả triplet của câu đó theo cú pháp~\eqref{eq:gas_format}.

Mỗi record có dạng:

\begin{equation}
  r = \bigl(\underbrace{x}_{\text{input\_text}},\;
            \underbrace{y}_{\text{target\_text}},\;
            \underbrace{\{a_i\}}_{\text{aspects}},\;
            \underbrace{\{(a_i, c_i, s_i)\}}_{\text{gold\_triples}}
      \bigr)
  \label{eq:gas_record}
\end{equation}

Tokenization được thực hiện bởi \texttt{T5Tokenizer} với
\texttt{max\_length}$= 128$ cho cả đầu vào lẫn đầu ra.  Các vị trí
padding trong chuỗi nhãn được thay bằng $-100$ để hàm mất mát
cross-entropy bỏ qua.

% ---------------------------------------------------------------------
\subsection{Huấn luyện}
\label{subsec:gas_train}
% ---------------------------------------------------------------------

Mô hình được tinh chỉnh bằng hàm mất mát cross-entropy chuẩn của
seq2seq (T5 tự tính theo cơ chế \textit{teacher forcing}):

\begin{equation}
  \mathcal{L}_{\text{GAS}} =
    -\sum_{t=1}^{T} \log P\!\left(y_t \mid y_{<t},\, \mathbf{x};\,
    \Theta\right)
  \label{eq:gas_loss}
\end{equation}

\noindent với $T$ là độ dài chuỗi đích và $\Theta$ là toàn bộ tham số
mô hình.  Quá trình tối ưu sử dụng AdamW với learning rate
$3 \times 10^{-4}$, kết hợp \textit{linear warmup} trong 10\% tổng
số bước.  Kỹ thuật mixed precision (AMP) được bật khi có GPU CUDA.

Early stopping theo dõi \textbf{ATE F1} (Aspect Term F1) trên tập
phát triển (patience$= 4$ epoch): một dự đoán được tính là True
Positive khi aspect term dự đoán khớp chính xác với nhãn vàng, bất
kể category hay sentiment — nhất quán với mục tiêu của Stage~2 trong
pipeline.

% ---------------------------------------------------------------------
\subsection{Hậu xử lý kết quả sinh}
\label{subsec:gas_postprocess}
% ---------------------------------------------------------------------

Chuỗi văn bản do T5 sinh ra được phân tích cú pháp bằng biểu thức
chính quy để trích xuất danh sách triplet:

\begin{equation}
  \texttt{\_GAS\_RE} =
    \texttt{\textbackslash(\textbackslash s*(.+?),\textbackslash s*%
    (.+?),\textbackslash s*(positive|negative|neutral)\textbackslash s*%
    \textbackslash)}
  \label{eq:gas_regex}
\end{equation}

Group~1, 2, 3 lần lượt trả về $a_i$, $c_i$, $s_i$.  Nếu chuỗi sinh
là \texttt{"none"} hoặc không khớp mẫu nào, hàm trả về danh sách
rỗng.

\paragraph{Chuẩn hóa aspect term bằng Levenshtein.}
Mô hình sinh đôi khi tạo ra cụm từ khía cạnh với lỗi chính tả nhỏ
(thiếu/thừa ký tự, hoán vị).  Để đảm bảo aspect term kết quả
\textit{thực sự xuất hiện trong câu gốc}, bước chuẩn hóa n-gram
Levenshtein được áp dụng sau khi parse:

\begin{enumerate}
  \item \textbf{Xây dựng từ điển n-gram:}  Từ câu gốc $\mathbf{x}$,
    tập ứng viên $\mathcal{V}$ gồm tất cả chuỗi con liên tục theo từ
    (unigram, bigram, trigram, \dots).

  \item \textbf{Ánh xạ về n-gram gần nhất:}  Nếu aspect term dự đoán
    $\hat{a}$ chưa thuộc $\mathcal{V}$, thay bằng:
    \begin{equation}
      a^* = \arg\min_{v \in \mathcal{V}}\;
        \mathrm{Lev}(\hat{a},\, v)
      \label{eq:levenshtein}
    \end{equation}
    \noindent trong đó $\mathrm{Lev}(\cdot, \cdot)$ là khoảng cách
    chỉnh sửa Levenshtein.  Nếu $\hat{a} \in \mathcal{V}$, giữ
    nguyên.
\end{enumerate}

Chỉ aspect term được chuẩn hóa; category và sentiment giữ nguyên
như mô hình sinh ra.

% ---------------------------------------------------------------------
\subsection{Độ đo đánh giá}
\label{subsec:gas_eval}
% ---------------------------------------------------------------------

Vì luận văn chỉ sử dụng GAS cho bước trích xuất aspect term (Stage~2),
độ đo duy nhất được báo cáo là \textbf{ATE F1} (Aspect Term F1):
một dự đoán được tính là True Positive khi chuỗi aspect term dự đoán
khớp chính xác (\textit{exact match}) với một aspect term trong nhãn vàng
của câu — category và sentiment không tham gia so khớp.

Precision, Recall và F1 được tính theo công thức chuẩn:

\begin{equation}
  P = \frac{\mathrm{TP}}{\mathrm{TP} + \mathrm{FP}}, \quad
  R = \frac{\mathrm{TP}}{\mathrm{TP} + \mathrm{FN}}, \quad
  F_1 = \frac{2PR}{P + R}
  \label{eq:prf}
\end{equation}

\noindent trong đó TP, FP, FN được tổng hợp (\textit{micro}) qua toàn
bộ câu của tập kiểm thử trước khi tính tỉ lệ — nhất quán với giao thức
đánh giá ATE trong các nghiên cứu ABSA~\cite{zhang2021gas,luo2019doer}.
 trong main.tex
% =====================================================================

% =====================================================================
\section{Phương pháp GAS (Generative Aspect Sentiment)}
\label{sec:gas}
% =====================================================================

GAS (\textit{Generative Aspect Sentiment}) là phương pháp trích xuất
ABSA một giai đoạn được đề xuất bởi Zhang et al.~\cite{zhang2021gas}.
Khác với các tiếp cận phân tầng truyền thống — đầu tiên trích xuất
aspect term (ATE), sau đó phân loại cảm xúc (APC) — GAS sử dụng một
mô hình sinh duy nhất để dự đoán đồng thời ba nhãn trong một lần giải
mã:
\textbf{aspect term}, \textbf{aspect category}, và
\textbf{sentiment polarity}.

Trong luận văn này, GAS được sử dụng làm \textit{baseline cho bước
trích xuất aspect term} (Stage~2 — ATE).  Mô hình vẫn được huấn luyện
để sinh bộ ba $(a_i, c_i, s_i)$ theo đúng cấu trúc gốc, song kết quả
so sánh chỉ tập trung vào nhãn \textbf{aspect term}: category và
sentiment không thuộc phạm vi đánh giá của chương này.

% ---------------------------------------------------------------------
\subsection{Kiến trúc tổng quan}
\label{subsec:gas_arch}
% ---------------------------------------------------------------------

GAS xây dựng bài toán ABSA như một bài toán dịch chuỗi
(\textit{sequence-to-sequence}), với mô hình nền là
\texttt{T5ForConditionalGeneration}~\cite{raffel2020t5} phiên bản
\texttt{t5-base}.  Với mỗi câu đầu vào $\mathbf{x}$, mô hình được
huấn luyện để sinh ra chuỗi đích $\mathbf{y}$ theo cú pháp cố định:

\begin{equation}
  \mathbf{y} = \texttt{(}a_1,\, c_1,\, s_1\texttt{);\;
               (}a_2,\, c_2,\, s_2\texttt{);\;} \dots
  \label{eq:gas_format}
\end{equation}

\noindent trong đó $a_i$ là cụm từ khía cạnh (\textit{aspect term}),
$c_i \in \{\text{SERVICE, FACILITY, AMENITY, EXPERIENCE, BRANDING,
LOYALTY}\}$ là danh mục khía cạnh, và
$s_i \in \{\text{positive, negative, neutral}\}$ là cực tính cảm xúc.
Nếu câu không chứa khía cạnh nào, mô hình sinh ra chuỗi
\texttt{"none"}.  Trong luận văn này chỉ trường $a_i$ được sử dụng
cho mục đích so sánh; $c_i$ và $s_i$ được Stage~3 (APC) xử lý riêng.

Quá trình giải mã sử dụng \textit{beam search} với chiều rộng tia
$k = 4$ và điều kiện dừng sớm (\textit{early stopping}) khi toàn bộ
các tia đều kết thúc bằng token \texttt{</s>}.

% ---------------------------------------------------------------------
\subsection{Tiền xử lý dữ liệu}
\label{subsec:gas_data}
% ---------------------------------------------------------------------

Dữ liệu đầu vào được đọc từ định dạng \texttt{.apc} bốn dòng lặp lại
(câu — aspect term — aspect category — sentiment).
Vì định dạng này lưu một dòng cho mỗi cặp \textit{(câu, aspect)},
các mẫu cùng thuộc một câu được nhóm lại thành một
\textit{record} duy nhất.  Chuỗi đích tương ứng được xây dựng bằng
cách nối tất cả triplet của câu đó theo cú pháp~\eqref{eq:gas_format}.

Mỗi record có dạng:

\begin{equation}
  r = \bigl(\underbrace{x}_{\text{input\_text}},\;
            \underbrace{y}_{\text{target\_text}},\;
            \underbrace{\{a_i\}}_{\text{aspects}},\;
            \underbrace{\{(a_i, c_i, s_i)\}}_{\text{gold\_triples}}
      \bigr)
  \label{eq:gas_record}
\end{equation}

Tokenization được thực hiện bởi \texttt{T5Tokenizer} với
\texttt{max\_length}$= 128$ cho cả đầu vào lẫn đầu ra.  Các vị trí
padding trong chuỗi nhãn được thay bằng $-100$ để hàm mất mát
cross-entropy bỏ qua.

% ---------------------------------------------------------------------
\subsection{Huấn luyện}
\label{subsec:gas_train}
% ---------------------------------------------------------------------

Mô hình được tinh chỉnh bằng hàm mất mát cross-entropy chuẩn của
seq2seq (T5 tự tính theo cơ chế \textit{teacher forcing}):

\begin{equation}
  \mathcal{L}_{\text{GAS}} =
    -\sum_{t=1}^{T} \log P\!\left(y_t \mid y_{<t},\, \mathbf{x};\,
    \Theta\right)
  \label{eq:gas_loss}
\end{equation}

\noindent với $T$ là độ dài chuỗi đích và $\Theta$ là toàn bộ tham số
mô hình.  Quá trình tối ưu sử dụng AdamW với learning rate
$3 \times 10^{-4}$, kết hợp \textit{linear warmup} trong 10\% tổng
số bước.  Kỹ thuật mixed precision (AMP) được bật khi có GPU CUDA.

Early stopping theo dõi \textbf{ATE F1} (Aspect Term F1) trên tập
phát triển (patience$= 4$ epoch): một dự đoán được tính là True
Positive khi aspect term dự đoán khớp chính xác với nhãn vàng, bất
kể category hay sentiment — nhất quán với mục tiêu của Stage~2 trong
pipeline.

% ---------------------------------------------------------------------
\subsection{Hậu xử lý kết quả sinh}
\label{subsec:gas_postprocess}
% ---------------------------------------------------------------------

Chuỗi văn bản do T5 sinh ra được phân tích cú pháp bằng biểu thức
chính quy để trích xuất danh sách triplet:

\begin{equation}
  \texttt{\_GAS\_RE} =
    \texttt{\textbackslash(\textbackslash s*(.+?),\textbackslash s*%
    (.+?),\textbackslash s*(positive|negative|neutral)\textbackslash s*%
    \textbackslash)}
  \label{eq:gas_regex}
\end{equation}

Group~1, 2, 3 lần lượt trả về $a_i$, $c_i$, $s_i$.  Nếu chuỗi sinh
là \texttt{"none"} hoặc không khớp mẫu nào, hàm trả về danh sách
rỗng.

\paragraph{Chuẩn hóa aspect term bằng Levenshtein.}
Mô hình sinh đôi khi tạo ra cụm từ khía cạnh với lỗi chính tả nhỏ
(thiếu/thừa ký tự, hoán vị).  Để đảm bảo aspect term kết quả
\textit{thực sự xuất hiện trong câu gốc}, bước chuẩn hóa n-gram
Levenshtein được áp dụng sau khi parse:

\begin{enumerate}
  \item \textbf{Xây dựng từ điển n-gram:}  Từ câu gốc $\mathbf{x}$,
    tập ứng viên $\mathcal{V}$ gồm tất cả chuỗi con liên tục theo từ
    (unigram, bigram, trigram, \dots).

  \item \textbf{Ánh xạ về n-gram gần nhất:}  Nếu aspect term dự đoán
    $\hat{a}$ chưa thuộc $\mathcal{V}$, thay bằng:
    \begin{equation}
      a^* = \arg\min_{v \in \mathcal{V}}\;
        \mathrm{Lev}(\hat{a},\, v)
      \label{eq:levenshtein}
    \end{equation}
    \noindent trong đó $\mathrm{Lev}(\cdot, \cdot)$ là khoảng cách
    chỉnh sửa Levenshtein.  Nếu $\hat{a} \in \mathcal{V}$, giữ
    nguyên.
\end{enumerate}

Chỉ aspect term được chuẩn hóa; category và sentiment giữ nguyên
như mô hình sinh ra.

% ---------------------------------------------------------------------
\subsection{Độ đo đánh giá}
\label{subsec:gas_eval}
% ---------------------------------------------------------------------

Vì luận văn chỉ sử dụng GAS cho bước trích xuất aspect term (Stage~2),
độ đo duy nhất được báo cáo là \textbf{ATE F1} (Aspect Term F1):
một dự đoán được tính là True Positive khi chuỗi aspect term dự đoán
khớp chính xác (\textit{exact match}) với một aspect term trong nhãn vàng
của câu — category và sentiment không tham gia so khớp.

Precision, Recall và F1 được tính theo công thức chuẩn:

\begin{equation}
  P = \frac{\mathrm{TP}}{\mathrm{TP} + \mathrm{FP}}, \quad
  R = \frac{\mathrm{TP}}{\mathrm{TP} + \mathrm{FN}}, \quad
  F_1 = \frac{2PR}{P + R}
  \label{eq:prf}
\end{equation}

\noindent trong đó TP, FP, FN được tổng hợp (\textit{micro}) qua toàn
bộ câu của tập kiểm thử trước khi tính tỉ lệ — nhất quán với giao thức
đánh giá ATE trong các nghiên cứu ABSA~\cite{zhang2021gas,luo2019doer}.
 trong main.tex
% =====================================================================

% =====================================================================
\section{Phương pháp GAS (Generative Aspect Sentiment)}
\label{sec:gas}
% =====================================================================

GAS (\textit{Generative Aspect Sentiment}) là phương pháp trích xuất
ABSA một giai đoạn được đề xuất bởi Zhang et al.~\cite{zhang2021gas}.
Khác với các tiếp cận phân tầng truyền thống — đầu tiên trích xuất
aspect term (ATE), sau đó phân loại cảm xúc (APC) — GAS sử dụng một
mô hình sinh duy nhất để dự đoán đồng thời ba nhãn trong một lần giải
mã:
\textbf{aspect term}, \textbf{aspect category}, và
\textbf{sentiment polarity}.

Trong luận văn này, GAS được sử dụng làm \textit{baseline cho bước
trích xuất aspect term} (Stage~2 — ATE).  Mô hình vẫn được huấn luyện
để sinh bộ ba $(a_i, c_i, s_i)$ theo đúng cấu trúc gốc, song kết quả
so sánh chỉ tập trung vào nhãn \textbf{aspect term}: category và
sentiment không thuộc phạm vi đánh giá của chương này.

% ---------------------------------------------------------------------
\subsection{Kiến trúc tổng quan}
\label{subsec:gas_arch}
% ---------------------------------------------------------------------

GAS xây dựng bài toán ABSA như một bài toán dịch chuỗi
(\textit{sequence-to-sequence}), với mô hình nền là
\texttt{T5ForConditionalGeneration}~\cite{raffel2020t5} phiên bản
\texttt{t5-base}.  Với mỗi câu đầu vào $\mathbf{x}$, mô hình được
huấn luyện để sinh ra chuỗi đích $\mathbf{y}$ theo cú pháp cố định:

\begin{equation}
  \mathbf{y} = \texttt{(}a_1,\, c_1,\, s_1\texttt{);\;
               (}a_2,\, c_2,\, s_2\texttt{);\;} \dots
  \label{eq:gas_format}
\end{equation}

\noindent trong đó $a_i$ là cụm từ khía cạnh (\textit{aspect term}),
$c_i \in \{\text{SERVICE, FACILITY, AMENITY, EXPERIENCE, BRANDING,
LOYALTY}\}$ là danh mục khía cạnh, và
$s_i \in \{\text{positive, negative, neutral}\}$ là cực tính cảm xúc.
Nếu câu không chứa khía cạnh nào, mô hình sinh ra chuỗi
\texttt{"none"}.  Trong luận văn này chỉ trường $a_i$ được sử dụng
cho mục đích so sánh; $c_i$ và $s_i$ được Stage~3 (APC) xử lý riêng.

Quá trình giải mã sử dụng \textit{beam search} với chiều rộng tia
$k = 4$ và điều kiện dừng sớm (\textit{early stopping}) khi toàn bộ
các tia đều kết thúc bằng token \texttt{</s>}.

% ---------------------------------------------------------------------
\subsection{Tiền xử lý dữ liệu}
\label{subsec:gas_data}
% ---------------------------------------------------------------------

Dữ liệu đầu vào được đọc từ định dạng \texttt{.apc} bốn dòng lặp lại
(câu — aspect term — aspect category — sentiment).
Vì định dạng này lưu một dòng cho mỗi cặp \textit{(câu, aspect)},
các mẫu cùng thuộc một câu được nhóm lại thành một
\textit{record} duy nhất.  Chuỗi đích tương ứng được xây dựng bằng
cách nối tất cả triplet của câu đó theo cú pháp~\eqref{eq:gas_format}.

Mỗi record có dạng:

\begin{equation}
  r = \bigl(\underbrace{x}_{\text{input\_text}},\;
            \underbrace{y}_{\text{target\_text}},\;
            \underbrace{\{a_i\}}_{\text{aspects}},\;
            \underbrace{\{(a_i, c_i, s_i)\}}_{\text{gold\_triples}}
      \bigr)
  \label{eq:gas_record}
\end{equation}

Tokenization được thực hiện bởi \texttt{T5Tokenizer} với
\texttt{max\_length}$= 128$ cho cả đầu vào lẫn đầu ra.  Các vị trí
padding trong chuỗi nhãn được thay bằng $-100$ để hàm mất mát
cross-entropy bỏ qua.

% ---------------------------------------------------------------------
\subsection{Huấn luyện}
\label{subsec:gas_train}
% ---------------------------------------------------------------------

Mô hình được tinh chỉnh bằng hàm mất mát cross-entropy chuẩn của
seq2seq (T5 tự tính theo cơ chế \textit{teacher forcing}):

\begin{equation}
  \mathcal{L}_{\text{GAS}} =
    -\sum_{t=1}^{T} \log P\!\left(y_t \mid y_{<t},\, \mathbf{x};\,
    \Theta\right)
  \label{eq:gas_loss}
\end{equation}

\noindent với $T$ là độ dài chuỗi đích và $\Theta$ là toàn bộ tham số
mô hình.  Quá trình tối ưu sử dụng AdamW với learning rate
$3 \times 10^{-4}$, kết hợp \textit{linear warmup} trong 10\% tổng
số bước.  Kỹ thuật mixed precision (AMP) được bật khi có GPU CUDA.

Early stopping theo dõi \textbf{ATE F1} (Aspect Term F1) trên tập
phát triển (patience$= 4$ epoch): một dự đoán được tính là True
Positive khi aspect term dự đoán khớp chính xác với nhãn vàng, bất
kể category hay sentiment — nhất quán với mục tiêu của Stage~2 trong
pipeline.

% ---------------------------------------------------------------------
\subsection{Hậu xử lý kết quả sinh}
\label{subsec:gas_postprocess}
% ---------------------------------------------------------------------

Chuỗi văn bản do T5 sinh ra được phân tích cú pháp bằng biểu thức
chính quy để trích xuất danh sách triplet:

\begin{equation}
  \texttt{\_GAS\_RE} =
    \texttt{\textbackslash(\textbackslash s*(.+?),\textbackslash s*%
    (.+?),\textbackslash s*(positive|negative|neutral)\textbackslash s*%
    \textbackslash)}
  \label{eq:gas_regex}
\end{equation}

Group~1, 2, 3 lần lượt trả về $a_i$, $c_i$, $s_i$.  Nếu chuỗi sinh
là \texttt{"none"} hoặc không khớp mẫu nào, hàm trả về danh sách
rỗng.

\paragraph{Chuẩn hóa aspect term bằng Levenshtein.}
Mô hình sinh đôi khi tạo ra cụm từ khía cạnh với lỗi chính tả nhỏ
(thiếu/thừa ký tự, hoán vị).  Để đảm bảo aspect term kết quả
\textit{thực sự xuất hiện trong câu gốc}, bước chuẩn hóa n-gram
Levenshtein được áp dụng sau khi parse:

\begin{enumerate}
  \item \textbf{Xây dựng từ điển n-gram:}  Từ câu gốc $\mathbf{x}$,
    tập ứng viên $\mathcal{V}$ gồm tất cả chuỗi con liên tục theo từ
    (unigram, bigram, trigram, \dots).

  \item \textbf{Ánh xạ về n-gram gần nhất:}  Nếu aspect term dự đoán
    $\hat{a}$ chưa thuộc $\mathcal{V}$, thay bằng:
    \begin{equation}
      a^* = \arg\min_{v \in \mathcal{V}}\;
        \mathrm{Lev}(\hat{a},\, v)
      \label{eq:levenshtein}
    \end{equation}
    \noindent trong đó $\mathrm{Lev}(\cdot, \cdot)$ là khoảng cách
    chỉnh sửa Levenshtein.  Nếu $\hat{a} \in \mathcal{V}$, giữ
    nguyên.
\end{enumerate}

Chỉ aspect term được chuẩn hóa; category và sentiment giữ nguyên
như mô hình sinh ra.

% ---------------------------------------------------------------------
\subsection{Độ đo đánh giá}
\label{subsec:gas_eval}
% ---------------------------------------------------------------------

Vì luận văn chỉ sử dụng GAS cho bước trích xuất aspect term (Stage~2),
độ đo duy nhất được báo cáo là \textbf{ATE F1} (Aspect Term F1):
một dự đoán được tính là True Positive khi chuỗi aspect term dự đoán
khớp chính xác (\textit{exact match}) với một aspect term trong nhãn vàng
của câu — category và sentiment không tham gia so khớp.

Precision, Recall và F1 được tính theo công thức chuẩn:

\begin{equation}
  P = \frac{\mathrm{TP}}{\mathrm{TP} + \mathrm{FP}}, \quad
  R = \frac{\mathrm{TP}}{\mathrm{TP} + \mathrm{FN}}, \quad
  F_1 = \frac{2PR}{P + R}
  \label{eq:prf}
\end{equation}

\noindent trong đó TP, FP, FN được tổng hợp (\textit{micro}) qua toàn
bộ câu của tập kiểm thử trước khi tính tỉ lệ — nhất quán với giao thức
đánh giá ATE trong các nghiên cứu ABSA~\cite{zhang2021gas,luo2019doer}.
